\scanpage{88}
\mspage{8}
% archive foliation: 87

schéma en carrés combinatoires sur $S$) dont le schéma des diagonales s'identifie à $\underline{\omega}_{\underline{J}}$ :
\[
(42)\qquad \underline{Q}_{E/S}/\unsure{\boldsymbol{\mu}|\mu}_2\simeq\underline{\omega}_{\underline{J}}
\]
Ceci posé, on voit que l'on a un isom\inserted{can.} de $(\unsure{\boldsymbol{\mu}|\mu}_4)_S$-torseurs
\[
(43)\qquad Q_{E/S}\wedge_{\unsure{\boldsymbol{\mu}|\mu}_4}Q_{0S}\simeq\mathcal{L}_{E/S}
\]
i.e.
\[
\mathcal{L}_{E/S}\simeq\underline{\operatorname{Isom}}_{\unsure{\boldsymbol{\mu}|\mu}_{4S}\text{-tors.}}(Q_{0S}^{\unsure{\wedge}-1},Q_{E/S})
\]
On aurait pu définir les épinglages de Legendre de $E/S$ comme les trivialisations de $Q_{E/S}$, au lieu de celles de $\mathcal{L}_{E/S}$ --- mais on aura alors une relation moins jolie avec les épinglages de Jacobi d'échelon $4$.

On voit en tous cas que les épinglages de Legendre\inserted{de $E/S$} forment\inserted{les sections d'} un torseur sous $\unsure{\boldsymbol{\mu}|\mu}_4$, tel\deleted{\unsure{le}} que le $\unsure{\boldsymbol{\mu}|\mu}_2$-torseur associé soit can. iso. à $\underline{\omega}_{\underline{J}}\inserted{\wedge\unsure{\boldsymbol{\mu}|\mu}_4^*}$
\[
(44)\quad\mathcal{L}_{E/S}/\unsure{\boldsymbol{\mu}|\mu}_2\simeq\mathcal{L}_{E/S}\wedge_{\unsure{\boldsymbol{\mu}|\mu}_{4S}}\unsure{\boldsymbol{\mu}|\mu}_{2S}\simeq\deleted{\underline{\omega}_{\underline{J}}}\underbrace{(\unsure{\boldsymbol{\mu}|\mu}_{4S}^*\wedge\underline{\omega}_{\underline{J}})}
\]
tandis que les épinglages de Jacobi d'échelon
